%-----先用xelatex编译，再用lualatex编译------

\documentclass[12pt,a4paper]{article}
\usepackage[UTF8]{ctex}
\usepackage{amsmath, amssymb, amsthm}
\usepackage{geometry}
\usepackage{enumitem}
\usepackage{fancyhdr}
\usepackage{titlesec}
\usepackage{xcolor}
\usepackage{fontspec}




%===========强制加粗，需用lualatex编译，在字体设置后加以下内容，若无需加粗可去掉=============

%		[
%	FakeBold=1.5,
%	AutoFakeBold=true
%	]



% ========== 中文字体，强制加粗====================

\setCJKmainfont{SimSun}[
BoldFont={SimHei},
FakeBold=1.5,
AutoFakeBold=true
]
\setCJKsansfont{SimHei}[
BoldFont={SimHei},
FakeBold=1.5,
AutoFakeBold=true
]
\setCJKfamilyfont{kai}{KaiTi}[
FakeBold=1.5,
AutoFakeBold=true
]
\setCJKfamilyfont{hei}{SimHei}[
FakeBold=1.5,
AutoFakeBold=true
]
\newcommand{\kai}{\CJKfamily{kai}}
\newcommand{\hei}{\CJKfamily{hei}}




% ========== 英文字体，强制加粗 ==========
% 数学模式使用 Times 风格
\usepackage{unicode-math}
\setmainfont{Times New Roman}[
FakeBold=1.5,
AutoFakeBold=true
]
\setmathfont{TeX Gyre Termes Math}[
FakeBold=1.5,
AutoFakeBold=true
]  % Times 风格的数学字体



% ========== 页面设置 ==========
\geometry{left=2cm, right=2cm, top=2.5cm, bottom=2.5cm}

% ========== 页眉页脚 ==========
\pagestyle{fancy}
\fancyhf{}
\cfoot{\small 高三数学 第\thepage 页（共4页）}
\renewcommand{\headrulewidth}{0pt}
\renewcommand{\footrulewidth}{0pt}



% ========== 选项环境 ==========
%1.选项较短用\choice{A}{B}{C}{D}
\newcommand{\choices}[4]{%
	\par\noindent
	\begin{tabular}{@{}p{0.22\linewidth}p{0.22\linewidth}p{0.22\linewidth}p{0.22\linewidth}@{}}
		A. #1 & B. #2 & C. #3 & D. #4
	\end{tabular}%
	\par
}

%2.选项较长用\Choice{A}{B}{C}{D}
\newcommand{\Choices}[4]{%
	\par\noindent
	\begin{tabular}{@{}p{0.44\linewidth}p{0.44\linewidth}@{}}
		A. #1 & B. #2
	\end{tabular}%
	\par
	\par\noindent
	\begin{tabular}{@{}p{0.44\linewidth}p{0.44\linewidth}@{}}
		C. #3 & D. #4
	\end{tabular}%
	\par
}


% ========== 列表间距 ==========
\setlist[enumerate,1]{label=\arabic*., left=0pt, itemsep=12pt, parsep=0pt, topsep=0pt}
\setlist[enumerate,2]{left=0pt, itemsep=2pt, labelsep=0.5em}

%enumerate 是 LaTeX 中用于创建自动编号列表的环境，可嵌套
%	\begin{enumerate}
%		\item 第一项
%		\item 第二项
%		\item 第三项
%	\end{enumerate}


%================================开始=======================================
\begin{document}

% ========== 第1页，需手动用\newpage 分页！==================
	% ========== 标题部分 ==========
	\begin{center}
		{\LARGE 清苑一中2026届高一年级六月供题}

		\vspace{1cm}
		{\hei\huge 数 \quad 学}
	\end{center}

	\vspace{0.3cm}

	% 双横线 + 左右两侧文字（楷体）
	\noindent
	\begin{minipage}[b]{0.5\textwidth}
		\raggedright
		{\kai 清苑一中高一数学组命制}
	\end{minipage}%
	\begin{minipage}[b]{0.5\textwidth}
		\raggedleft
		{\kai 2026.6}
	\end{minipage}
	\\[-11pt]
	\noindent\rule{\textwidth}{1pt}\\[-16pt]
	\noindent\rule{\textwidth}{1pt}
	\\[-20pt]

	% ========== 试卷说明 ==========
	{\hei\noindent\hspace{2em}本卷共4页，19题，全卷满分150分。用时120分钟。}

	\vspace{0.5cm}

	% ========== 注意事项 ==========
	{\hei\noindent 注意事项：}
	\begin{enumerate}[left=2em, itemsep=2pt, labelsep=0.5em, topsep=0pt]
		\item {\kai 选择题的作答：每小题选出答案后，用2B铅笔把答题卡上对应题目的答案标号涂黑。}
		\item {\kai 非选择题的作答：用黑色签字笔直接答在答题卡上对应的答题区域内。}
	\end{enumerate}

	\vspace{0.3cm}

	% ========== 单选 =========
	{\hei\noindent 一、单项选择题：本题共8小题，每小题5分，共40分。在每小题给出的四个选项中，只有一项是符合题目要求。}

	\vspace{0.3cm}

	\begin{enumerate}
		\item 若 \( a, b \) 为实数，且 \( \dfrac{a}{2+i} = b-i \)，则 \( a+b= \)
		\choices{7}{5}{-5}{-7}

		\item
		\choices{}{}{}{}

		\item
		\choices{}{}{}{}

		\item
		\choices{}{}{}{}

		\item
		\choices{}{}{}{}
	\end{enumerate}

	\newpage

% ========== 第2页 ==========================================================


	%============多选==============
	{\hei\noindent 二、多项选择题：本题共3小题，每小题6分，共18分。在每小题给出的选项中，有多项符合题目要求。全部选对的得6分，部分选对的得部分分，有选错的得0分。}

	\vspace{0.3cm}

	\begin{enumerate}[resume]	%[resume] 是 enumitem 宏包提供的一个选项，用于让编号从上一次结束的地方继续。


		\item 在 \( (2x - \dfrac{1}{x})^2 (x^2 + 2) \) 的展开式中，含 \( x^2 \) 项的系数为
		\choices{240}{-240}{80}{-80}

		\item
		\choices{}{}{}{}

		\item
		\choices{}{}{}{}
	\end{enumerate}

	\vspace{0.5cm}



	%===================填空=======================================
	{\hei\noindent 三、填空题：本题共3小题，每小题5分，共15分。}

	\vspace{0.3cm}

	\begin{enumerate}[resume]
		\item 已知向量 \( \vec{a} = (1, 2) \)，\( \vec{b} = (3, -4) \)，则 \( \vec{a} \cdot \vec{b} = \)\underline{\hspace{3cm}}

		\item
		\underline{\hspace{3cm}}

		\item
		\underline{\hspace{3cm}}
	\end{enumerate}

	\vspace{0.5cm}


	%=========================大题===========================================
	{\hei\noindent 四、解答题：本题共5小题，共77分。解答应写出文字说明、证明过程或演算步骤。}

	\vspace{0.3cm}

	\begin{enumerate}[resume]
		\item (13分) 已知数列 \( \{a_n\} \) 的前 \( n \) 项和为 \( S_n \)，且 \( S_n = 2a_n - 1 \)。
		\begin{enumerate}
			\item 求数列 \( \{a_n\} \) 的通项公式；
			\item 设 \( b_n = \log_2 a_n \)，求数列 \( \{b_n\} \) 的前 \( n \) 项和 \( T_n \)。
		\end{enumerate}



		\item (  分)
		\begin{enumerate}
			\item
			\item
		\end{enumerate}



		\item ( 分)
		\begin{enumerate}
			\item
			\item
		\end{enumerate}



		\item ( 分)
		\begin{enumerate}
			\item
			\item
		\end{enumerate}



		\item ( 分)
		\begin{enumerate}
			\item
			\item
			\item
		\end{enumerate}


	\end{enumerate}




\end{document}
